\section*{Chapter \ref{ch:g10:synthesis-yield} --- Synthesis: Yield and Purity}

\begin{solution}{exo:g10:synthesis-yield:1}
$\eta = 4.2 / 5.0 = 0.84$, a \omterm{def:g10:synthesis-yield:yield}{yield} of \qty{84}{\percent}.
\end{solution}

\begin{solution}{exo:g10:synthesis-yield:2}
Weighing the \omterm{def:g7:chemical-reactions:reactant}{reactants}; \omterm{def:g10:synthesis-yield:reflux}{heating under reflux}; \omterm{def:g10:synthesis-yield:vacuum-filtration}{vacuum filtration} of the
\omterm{def:g10:synthesis-yield:synthesis}{crude product}; \omterm{def:g10:synthesis-yield:recrystallisation}{recrystallisation}; measuring the melting point.
\end{solution}

\begin{solution}{exo:g10:synthesis-yield:3}
Fed from the bottom, the water fills the whole jacket before leaving at
the top, so the whole length of the condenser is cooled. The top is left
open so that the apparatus is not closed: heating a closed apparatus
would build up pressure and could make it burst.
\end{solution}

\begin{solution}{exo:g10:synthesis-yield:4}
The product must be very soluble in the hot \omterm{def:g6:solutions-and-solubility:solute}{solvent} and hardly soluble
in the cold \omterm{def:g6:solutions-and-solubility:solute}{solvent}. The impurities, present in small amounts, stay
dissolved in the cold \omterm{def:g6:solutions-and-solubility:solute}{solvent} and end up in the \omterm{def:g3:separating-mixtures:filter}{filtrate}.
\end{solution}

\begin{solution}{exo:g10:synthesis-yield:5}
No: it melts below \qty{135}{\celsius} and over a range of several
degrees, the sign of an impure solid. It should be recrystallised, dried,
and its melting point measured again.
\end{solution}

\begin{solution}{exo:g10:synthesis-yield:6}
$n = 2.00 / 138.0 = \qty{0.0145}{mol}$ of salicylic acid, the \omterm{def:g10:reaction-progress-table:limiting}{limiting
reactant}; $m_{\max} = 0.0145 \times 180.0 = \qty{2.61}{g}$;
$\eta = 2.03 / 2.61 = 0.78$, that is \qty{78}{\percent}.
\end{solution}

\begin{solution}{exo:g10:synthesis-yield:7}
The weighed mass includes water, so it is larger than the mass of
aspirin: the computed \omterm{def:g10:synthesis-yield:yield}{yield} is too high, here impossibly above
\qty{100}{\percent}. The dry \omterm{def:g10:synthesis-yield:synthesis}{crude product} still contains impurities
(unreacted salicylic acid, by-products), which add mass too: its \omterm{def:g10:synthesis-yield:yield}{yield}
is also too high compared with the true \omterm{def:g10:synthesis-yield:yield}{yield} of pure aspirin.
\end{solution}

\begin{solution}{exo:g10:synthesis-yield:8}
It is much faster, and it leaves the crystals much drier, since the
suction draws out most of the liquid; the crystals are also easy to wash
on the funnel and to scrape off the flat filter paper.
\end{solution}

\begin{solution}{exo:g10:synthesis-yield:9}
Product $3.1/5.0 = 0.62$; aspirin $3.1/5.0 = 0.62$; salicylic acid
$2.2/5.0 = 0.44$. The product travels like aspirin and shows no spot of
salicylic acid: it is aspirin, with no salicylic acid detected.
\end{solution}

\begin{solution}{exo:g10:synthesis-yield:10}
\omterm{def:g6:solutions-and-solubility:solute}{Solvent} A: very soluble hot, hardly soluble cold; \omterm{def:g6:solutions-and-solubility:solute}{solvent} B keeps most of
the product dissolved even cold. With A, \qty{3.0}{g} need about
$3.0 / 150 = \qty{0.020}{L}$, that is \qty{20}{mL} of hot \omterm{def:g6:solutions-and-solubility:solute}{solvent}; once
cold, at most $2 \times 0.020 = \qty{0.04}{g}$ stays dissolved.
\end{solution}

\begin{solution}{exo:g10:synthesis-yield:11}
A boiling liquid needs room above it: froth and splashes must not reach
the condenser. Boiling stones give the vapour small places to form
bubbles, so that the liquid boils smoothly instead of in sudden bursts.
\end{solution}

\begin{solution}{exo:g10:synthesis-yield:12}
$0.80 \times 0.70 = 0.56$: \qty{56}{\percent}. Three steps of
\qty{90}{\percent}: $0.90^3 = 0.73$, \qty{73}{\percent}. Every step
loses product (and time and money), and the losses multiply: the fewer
steps, the more product is left at the end.
\end{solution}

\begin{solution}{exo:g10:synthesis-yield:13}
At most $3.3 \times 0.050 = \qty{0.17}{g}$ of aspirin. Ice-cold water
\omterm{def:g2:mixing-and-dissolving:dissolve}{dissolves} even less aspirin, since its \omterm{def:g6:solutions-and-solubility:solubility}{solubility} falls with the
temperature.
\end{solution}

\begin{solution}{exo:g10:synthesis-yield:14}
\begin{enumerate}
  \item $n_0(\text{salicylic acid}) = 2.76 / 138.0 = \qty{0.0200}{mol}$,
        anhydride \qty{0.0500}{mol}: the salicylic acid is limiting,
        $x_{\max} = \qty{0.0200}{mol}$, and $0.0500 - 0.0200 =
        \qty{0.0300}{mol}$ of anhydride is left.
  \item The \omterm{def:g10:synthesis-yield:synthesis}{synthesis} gives \qty{0.0200}{mol} of ethanoic acid, the
        destruction of the excess $2 \times 0.0300 = \qty{0.0600}{mol}$:
        \qty{0.0800}{mol} in all.
\end{enumerate}
\end{solution}

\begin{solution}{exo:g10:synthesis-yield:15}
$n(\text{aspirin}) = \num{1.00e6} / 180.0 = \qty{5.56e3}{mol}$. With a
\omterm{def:g10:synthesis-yield:yield}{yield} of \qty{85}{\percent}, salicylic acid needed:
$\num{5.56e3} / 0.85 = \qty{6.54e3}{mol}$, that is
$\num{6.54e3} \times 138.0 = \qty{9.0e5}{g}$, about \qty{0.90}{t}.
\end{solution}

\begin{solution}{pb:g10:synthesis-yield:1}
\textbf{1.} C: $7 + 4 = 11$ and $9 + 2 = 11$; H: $6 + 6 = 12$ and
$8 + 4 = 12$; O: $3 + 3 = 6$ and $4 + 2 = 6$. Balanced.

\textbf{2.} Salicylic acid $7 \times 12.0 + 6 \times 1.0 + 3 \times 16.0
= \qty{138.0}{g/mol}$; anhydride $4 \times 12.0 + 6 \times 1.0 +
3 \times 16.0 = \qty{102.0}{g/mol}$; aspirin $9 \times 12.0 + 8 \times
1.0 + 4 \times 16.0 = \qty{180.0}{g/mol}$; ethanoic acid $2 \times 12.0
+ 4 \times 1.0 + 2 \times 16.0 = \qty{60.0}{g/mol}$.

\textbf{3.} \omterm{def:g7:chemical-reactions:reactant}{Reactants}: salicylic acid and ethanoic anhydride; product
wanted: aspirin; by-product: ethanoic acid.

\textbf{4.} Heating makes the reaction faster; under reflux the vapours
are condensed and returned, so no \omterm{def:g7:chemical-reactions:reactant}{reactant} or product is lost.

\textbf{5.} $n = 3.0 / 138.0 = \qty{0.0217}{mol}$.

\textbf{6.} $m = 6.0 \times 1.08 = \qty{6.48}{g}$;
$n = 6.48 / 102.0 = \qty{0.0635}{mol}$.

\textbf{7.} Salicylic acid: $0.0217 - x$; anhydride $0.0635 - x$;
aspirin and ethanoic acid: $x$. Salicylic acid runs out first:
$x_{\max} = \qty{0.0217}{mol}$, salicylic acid is limiting.

\textbf{8.} $m_{\max} = 0.0217 \times 180.0 = \qty{3.91}{g}$.

\textbf{9.} So that all the salicylic acid is turned into aspirin: any
left over would stay mixed with the aspirin and be hard to remove,
whereas the excess anhydride is destroyed by the water added at the end,
and the ethanoic acid it gives is washed away.

\textbf{10.} It includes water: \qty{3.6}{g} would give a false
$3.6 / 3.91 = \qty{92}{\percent}$.

\textbf{11.} $3.3 / 3.91 = 0.84$: \qty{84}{\percent} of crude dry
product, which is not pure.

\textbf{12.} The impurities are removed, and part of the aspirin stays
dissolved in the cold \omterm{def:g6:solutions-and-solubility:solute}{solvent} and in the washings.

\textbf{13.} $3.3 \times 0.030 = \qty{0.10}{g}$ at most.

\textbf{14.} $\eta = 2.7 / 3.91 = 0.69$: \qty{69}{\percent}.

\textbf{15.} A sharp melting point at the value of pure aspirin: the
crystals are aspirin, and pure.

\textbf{16.} Melting below \qty{135}{\celsius} and over a wide range:
the \omterm{def:g10:synthesis-yield:synthesis}{crude product} was impure.

\textbf{17.} The \omterm{def:g10:synthesis-yield:synthesis}{crude product} contained aspirin and unreacted salicylic
acid; the \omterm{def:g10:synthesis-yield:recrystallisation}{recrystallisation} removed the salicylic acid.

\textbf{18.} $R_f = 3.1 / 5.0 = 0.62$.

\textbf{19.} About as much in each: from \qty{3.91}{g} possible to
\qty{3.3}{g} of \omterm{def:g10:synthesis-yield:synthesis}{crude product}, and from \qty{3.3}{g} to \qty{2.7}{g} in
the \omterm{def:g10:synthesis-yield:recrystallisation}{recrystallisation}, about \qty{0.6}{g} each. Since the \omterm{def:g10:synthesis-yield:synthesis}{crude product}
was not pure aspirin, the first part of the work in fact lost a little
more aspirin than \qty{0.6}{g}.

\textbf{20.} The \omterm{def:g10:synthesis-yield:yield}{yield} of the \omterm{def:g10:synthesis-yield:synthesis}{synthesis} is about \qty{69}{\percent}.
\end{solution}
